\section{Metode}

\subsection{Målekjede}
Oppsettet følger en standard målekjede for tilstandsovervåking \parencite{tbv2}: vibrasjonen fanges opp av et akselerometer, sendes til en datainnsamlingsenhet (SKF IMx-8), og presenteres i analyseprogramvare (SKF Observer) der frekvensanalysen gjøres.

\subsection{Beskrivelse av oppstilling}
Forsøket ble gjennomført på en RK4-vibrasjonsrigg fra Bently Nevada. Riggen består av en DC-motor som driver en aksel opplagret i lagerbukker. Turtallet ble kontrollert og målt underveis, og vi kjørte målinger ved omtrent 2500, 4000 og 6000 rpm.

Vibrasjonene ble målt med piezoelektriske akselerometre fra SKF montert ved lageret \parencite{vibrasjon1}. Sensorene var koblet til ett SKF IMx-8 system, og signalene ble analysert i SKF Observer der vi leste av frekvenser og amplituder fra frekvensspekteret.

Sensorene ble plassert nær lageret for å redusere demping mellom lager og målepunkt \parencite{vibrasjon3}.

\subsubsection*{Lager brukt i forsøket}
I forsøket ble det brukt et R10-kulelager fra National Precision Bearing. Dette er et enradig kulelager med kontaktvinkel $\theta \approx 0°$ \parencite{beregninglager}.
\begin{table}[h!]
\centering
\begin{tabular}{|l|c|}
\hline
\textbf{Lagertype} & \textbf{Åpent rulle-/kulelager (R10)} \\ \hline
Antall kuler & 10 \\ \hline
Pitch diameter (mm) & 25.4 \\ \hline
Kulediameter (mm) & 4.7 \\ \hline
Indre diameter (mm) & 15.9 \\ \hline
Ytre diameter (mm) & 34.9 \\ \hline
\end{tabular}
\caption{Lagerdata for R10}
\end{table}

% TODO: Beskriv fremgangsmåte steg-for-steg (hvilke turtall, hvordan data ble hentet)
% TODO: Legg inn bilde av oppstillingen (fra øvingsarket eller eget foto)
